\documentclass[12pt]{article}
\usepackage[utf8]{inputenc}
\usepackage[T1]{fontenc}
\usepackage{hyperref}
\title{Banda-Yangere Language Resources for K-12 Curriculum Development}
\author{Compiled by Claude Sonnet 5.5 (claude-sonnet-5-5) for Andr\'as Kornai}
\date{}
\begin{document}
\maketitle

\section{Introduction}
Banda-Yangere (Yangere, ISO 639-3 \texttt{yaj}) is a Ubangian language of the Central African Republic, spoken around Berberati and closely related to Central Banda. Sources found in search give about 27,000 speakers, from an older count. Wikitongues lists the language as unwritten. The search found no school materials, so this report is short.

\section{Primary Educational Resources}
No primers, textbooks or curricula for Banda-Yangere turned up in the search.

\section{Linguistic and Reference Materials}
\begin{itemize}
\item Wikipedia, Yangere language (overview and classification): \url{https://en.wikipedia.org/wiki/Yangere_language}
\item Wikitongues language page for \texttt{yaj}: \url{https://wikitongues.org/languages/yaj/}
\end{itemize}
No grammar or dictionary was found in the results.

\section{Religious and Community Resources}
\begin{itemize}
\item Joshua Project, Banda-Yangere in Central African Republic (people-group data): \url{https://joshuaproject.net/people_groups/15981}
\end{itemize}

\section{Digital and Audio Resources}
\begin{itemize}
\item Global Recordings Network, language page for \texttt{yaj} (audio recordings listed): \url{https://globalrecordings.net/en/language/yaj}
\end{itemize}

\section{Curriculum Development Context}
The language is unwritten and no teaching materials were found. A curriculum would first need an orthography, community consultation and documentation. The Banda literature, including Moñino's work cited by Wikipedia, may help, but I did not see it directly.
\end{document}